% Bagian: first-order-logic
% Bab: syntax-and-semantics
% Subbagian: main-operator

\documentclass[../../../include/open-logic-section]{subfiles}

\begin{document}

\olfileid[id]{fol}{syn}{mai}

\olsection{\printtoken{S}{main operator} Formula}

\begin{explain}
Sering kali berguna untuk membicarakan operator terakhir yang digunakan
dalam membangun !!a{formula}~$!A$. Operator ini disebut \emph{!!{main
  operator}} dari~$!A$. Secara intuitif, operator tersebut merupakan
operator ``terluar'' dari $!A$. Misalnya, !!{main operator} dari $\lnot !A$
ialah $\lnot$, !!{main operator} dari $(!A \lor !B)$ ialah $\lor$, dan
seterusnya.
\end{explain}


\begin{defn}[!!^{main operator}]
\ollabel{def:main-op}
\emph{!!{main operator}} suatu !!{formula}~$!A$ didefinisikan sebagai
berikut:
\begin{enumerate}
\item \indcase*{!A}{!A}{$\indfrm$ tidak memiliki !!{main operator}.}

\tagitem{prvNot}{\indcase{!A}{\lnot !B}{!!{main operator} dari $\indfrm$
  ialah~$\lnot$.}}{}

\tagitem{prvAnd}{\indcase{!A}{(!B \land !C)}{!!{main operator} dari
  $\indfrm$ ialah~$\land$.}}{}

\tagitem{prvOr}{\indcase{!A}{(!B \lor !C)}{!!{main operator} dari
  $\indfrm$ ialah~$\lor$.}}{}

\tagitem{prvIf}{\indcase{!A}{(!B \lif !C)}{!!{main operator} dari
  $\indfrm$ ialah~$\lif$.}}{}

\tagitem{prvIff}{\indcase{!A}{(!B \liff !C)}{!!{main operator} dari
  $\indfrm$ ialah~$\liff$.}}{}

\tagitem{prvAll}{\indcase{!A}{\lforall[x][!B]}{!!{main operator} dari
  $\indfrm$ ialah~$\lforall$.}}{}

\tagitem{prvEx}{\indcase{!A}{\lexists[x][!B]}{!!{main operator} dari
  $\indfrm$ ialah~$\lexists$.}}{}
\end{enumerate}
\end{defn}

Dalam setiap kasus, yang kita maksud ialah \emph{kemunculan} tertentu dari
!!{main operator} yang ditunjukkan dalam formula. Misalnya, karena formula
$((!D \lif !E) \lif (!E \lif !D))$ berbentuk $(!B \lif !C)$ dengan $!B$
ialah $(!D \lif !E)$ dan $!C$ ialah $(!E \lif !D)$, kemunculan kedua
$\lif$ merupakan !!{main operator}.

\begin{explain}
Ini merupakan definisi \emph{rekursif} suatu fungsi yang memetakan semua
!!{formula}s nonatomik ke kemunculan !!{main operator}-nya. Karena cara
!!{formula}s didefinisikan secara induktif, setiap !!{formula}~$!A$
memenuhi salah satu kasus dalam \olref{def:main-op}. Hal ini menjamin bahwa
untuk setiap !!{formula} nonatomik~$!A$, terdapat suatu !!{main operator}.
Karena setiap !!{formula} hanya memenuhi satu dari kondisi-kondisi tersebut,
dan karena !!{formula}s yang lebih kecil yang digunakan untuk membangun
$!A$ ditentukan secara unik dalam setiap kasus, kemunculan !!{main
  operator} dari~$!A$ juga unik. Dengan demikian, kita telah mendefinisikan
suatu fungsi.
\end{explain}

Kita menamai !!{formula}s menurut \olref{tab:main-op}, bergantung pada simbol
yang menjadi !!{main operator}-nya.
\iftag{defNot,defOr,defAnd,defIf,defIff,defTrue,defFalse,defEx,defAll}
{ Namun, ingatlah bahwa operator terdefinisi secara resmi tidak muncul dalam
!!{formula}s. Operator tersebut hanyalah singkatan, sehingga secara resmi
tidak dapat menjadi operator utama suatu formula. Oleh karena itu, dalam
bukti tentang semua !!{formula}s, operator tersebut tidak perlu ditangani
secara tersendiri.}

\begin{table}[!h]
\centering
\begin{tabular}{c | c | c}
!!^{main operator} & Jenis !!{formula} & Contoh\\
\hline
tidak ada & atomik (!!{formula}) &
\iftag{prvFalse}{$\lfalse$,}{}
\iftag{prvTrue}{$\ltrue$,}{}
$\Atom{R}{t_1, \dots, t_n}$,
$\eq[t_1][t_2]$\\
$\lnot$ & negasi & $\lnot !A$ \\
$\land$ & konjungsi & $(!A \land !B)$ \\
$\lor$ & disjungsi & $(!A \lor !B)$ \\
$\lif$ & !!{conditional} & $(!A \lif !B)$ \\
$\liff$ & !!{biconditional} & $(!A \liff !B)$ \\
$\lforall[][]$ & !!^{formula} universal & $\lforall[x][!A]$ \\
$\lexists[][]$ & !!^{formula} eksistensial & $\lexists[x][!A]$
\end{tabular}
\caption{Operator utama dan nama !!{formula}s}
\ollabel{tab:main-op}
\end{table}

\end{document}
